%==============================================================
%  lecons/algebre/hyperplans_matrices.tex
%  Hyperplans de M_n(R) : intersection avec GL_n et O_n
%==============================================================

\lecontitle{Hyperplans de $\mathcal{M}_n(\mathbb{R})$}{Algèbre linéaire — Rencontre avec $\mathrm{GL}_n(\mathbb{R})$ et $\mathrm{O}_n(\mathbb{R})$}

%=============================================================
%  § 0 — PRÉLIMINAIRE : FORMES LINÉAIRES ET TRACE
%=============================================================
\secnum{navyblue}{1}{Préliminaire : hyperplans et représentation par la trace}

\begin{tcolorbox}[thmbox]
  \textbf{Cadre.}\enspace%
  \index{hyperplan de matrices}\index{forme linéaire}
  On travaille dans l'espace $\mathcal{M}_n(\mathbb{R})$ des matrices carrées
  réelles, de dimension $n^2$. Un \emph{hyperplan} $H$ est un sous-espace de
  dimension $n^2 - 1$, c'est-à-dire le noyau d'une forme linéaire non nulle.

  \medskip
  \textbf{Théorème (représentation des formes linéaires).}\enspace%
  L'application
  \[
    \Phi : \mathcal{M}_n(\mathbb{R}) \longrightarrow \mathcal{M}_n(\mathbb{R})^*,
    \qquad A \longmapsto \bigl(M \mapsto \mathrm{tr}(AM)\bigr)
  \]
  est un \emph{isomorphisme}. Autrement dit, toute forme linéaire
  $\varphi$ sur $\mathcal{M}_n(\mathbb{R})$ s'écrit de manière unique
  \[
    \varphi(M) = \mathrm{tr}(AM), \qquad A \in \mathcal{M}_n(\mathbb{R}).
  \]
  Tout hyperplan est donc de la forme
  \[
    H_A = \{\,M \in \mathcal{M}_n(\mathbb{R}) \mid \mathrm{tr}(AM) = 0\,\},
    \qquad A \neq 0.
  \]
\end{tcolorbox}

\begin{mdframed}[style=proofbar]
  $\Phi$ est linéaire et part d'un espace de dimension $n^2$ vers son dual,
  de même dimension : il suffit de vérifier l'injectivité. Si $\mathrm{tr}(AM)=0$
  pour toute $M$, on choisit $M = E_{ji}$ (matrice élémentaire) et l'on obtient
  $\mathrm{tr}(AE_{ji}) = A_{ij} = 0$ pour tous $i,j$, donc $A = 0$. \hfill$\blacksquare$
\end{mdframed}

\rubriq{Question}
Un hyperplan $H_A$ est un objet « générique » de codimension $1$. Contient-il
nécessairement des matrices remarquables — \emph{inversibles}, voire
\emph{orthogonales} ? La réponse est \textbf{oui} dès que $n \geq 2$, et c'est
l'objet des deux parties qui suivent. (Pour $n=1$, $\mathcal{M}_1(\mathbb{R})=\mathbb{R}$
et le seul hyperplan est $\{0\}$, qui ne rencontre ni $\mathrm{GL}_1$ ni $\mathrm{O}_1$.)

\decsep{}

%=============================================================
%  PARTIE I — GL_n(R)
%=============================================================
\secnum{navyblue}{2}{Partie I — Tout hyperplan rencontre $\mathrm{GL}_n(\mathbb{R})$}

\begin{tcolorbox}[thmbox]
  \textbf{Théorème.}\enspace%
  \index{groupe linéaire}
  Soit $n \geq 2$. Tout hyperplan $H$ de $\mathcal{M}_n(\mathbb{R})$ contient
  une matrice inversible :
  \[
    H \cap \mathrm{GL}_n(\mathbb{R}) \neq \varnothing.
  \]
\end{tcolorbox}

\rubriq{Idée}
On écrit $H = \ker\bigl(M\mapsto\mathrm{tr}(AM)\bigr)$ avec $A\neq 0$, puis on
ramène $A$ à une matrice $J_r = \mathrm{diag}(I_r,0)$ par équivalence. Il reste
à exhiber une matrice inversible dont les $r$ premiers coefficients diagonaux
ont une somme nulle : la \emph{matrice de permutation cyclique} convient, car
elle a une diagonale entièrement nulle.

\begin{mdframed}[style=proofbar]
  \textit{Réduction.}\enspace%
  Soit $A \neq 0$ avec $H = \{M : \mathrm{tr}(AM)=0\}$. Posons $r = \mathrm{rg}(A) \geq 1$.
  Il existe $P,Q \in \mathrm{GL}_n(\mathbb{R})$ telles que $A = P\,J_r\,Q$, où
  $J_r = \mathrm{diag}(I_r,\,0_{n-r})$. Pour toute $M$,
  \[
    \mathrm{tr}(AM) = \mathrm{tr}(P J_r Q M) = \mathrm{tr}(J_r\,QMP).
  \]
  L'application $M \mapsto N := QMP$ est une bijection de $\mathrm{GL}_n(\mathbb{R})$
  sur lui-même. Il suffit donc de trouver $N \in \mathrm{GL}_n(\mathbb{R})$ avec
  \[
    \mathrm{tr}(J_r N) = \sum_{i=1}^{r} N_{ii} = 0,
  \]
  car alors $M = Q^{-1}N P^{-1} \in H \cap \mathrm{GL}_n(\mathbb{R})$.

  \medskip
  \textit{Construction.}\enspace%
  \index{matrice de permutation}
  Soit $C$ la matrice de la permutation cyclique $\gamma = (1\,2\,\cdots\,n)$,
  définie par $C e_i = e_{i+1}$ (indices modulo $n$). Comme $n \geq 2$, on a
  $\gamma(i) \neq i$ pour tout $i$, donc \emph{tous} les coefficients diagonaux
  de $C$ sont nuls : en particulier $\sum_{i=1}^{r} C_{ii} = 0$. De plus $C$ est
  une matrice de permutation, donc inversible ($\det C = (-1)^{n-1}$). Le choix
  $N = C$ conclut. \hfill$\blacksquare$
\end{mdframed}

\rubriq{Remarque}
L'énoncé voisin « $\mathrm{GL}_n(\mathbb{R})$ n'est contenu dans aucun
hyperplan » est différent (et plus facile) : il résulte de $\mathrm{Vect}\,\mathrm{GL}_n(\mathbb{R}) = \mathcal{M}_n(\mathbb{R})$.
Ici on affirme bien plus : \emph{chaque} hyperplan, si fin soit-il, attrape une
matrice inversible.

\decsep{}

%=============================================================
%  PARTIE II — O_n(R)
%=============================================================
\secnum{navyblue}{3}{Partie II — Extension : tout hyperplan rencontre $\mathrm{O}_n(\mathbb{R})$}

\begin{tcolorbox}[thmbox]
  \textbf{Théorème.}\enspace%
  \index{groupe orthogonal}
  Soit $n \geq 2$. Tout hyperplan $H$ de $\mathcal{M}_n(\mathbb{R})$ contient
  une matrice orthogonale :
  \[
    H \cap \mathrm{O}_n(\mathbb{R}) \neq \varnothing.
  \]
\end{tcolorbox}

\rubriq{Idée}
On écrit $H = \ker\bigl(M\mapsto\mathrm{tr}(AQ)\bigr)$. La
\emph{décomposition en valeurs singulières}\index{décomposition en valeurs singulières}
réduit $A$ à une matrice diagonale $S = \mathrm{diag}(\sigma_1,\ldots,\sigma_n)$
à coefficients $\geq 0$, sans modifier le groupe orthogonal sur lequel on
optimise. On considère alors le \emph{chemin continu} de rotations $\theta \mapsto R(\theta)$
et l'on applique le théorème des valeurs intermédiaires à
$g(\theta) = \mathrm{tr}\bigl(S\,R(\theta)\bigr)$.

\begin{mdframed}[style=proofbar]
  \textit{Réduction orthogonale (SVD).}\enspace%
  Soit $A \neq 0$ avec $H = \{Q : \mathrm{tr}(AQ)=0\}$. La décomposition en
  valeurs singulières fournit $U,V \in \mathrm{O}_n(\mathbb{R})$ et
  $S = \mathrm{diag}(\sigma_1,\ldots,\sigma_n)$ avec $\sigma_i \geq 0$, telles
  que $A = U S V$. Pour $Q \in \mathrm{O}_n(\mathbb{R})$,
  \[
    \mathrm{tr}(AQ) = \mathrm{tr}(U S V Q) = \mathrm{tr}\bigl(S\,(VQU)\bigr).
  \]
  Quand $Q$ parcourt $\mathrm{O}_n(\mathbb{R})$, $R := VQU$ parcourt aussi
  $\mathrm{O}_n(\mathbb{R})$. Il suffit donc de trouver $R \in \mathrm{O}_n(\mathbb{R})$
  avec
  \[
    \mathrm{tr}(SR) = \sum_{i=1}^{n} \sigma_i\,R_{ii} = 0,
  \]
  car alors $Q = V^{\!\top} R\, U^{\!\top} \in H \cap \mathrm{O}_n(\mathbb{R})$.
  Quitte à conjuguer $S$ par une matrice de permutation (orthogonale), on
  suppose les $\sigma_i$ rangés de sorte que $\sigma_n = \min_i \sigma_i$.
  Notons $\tau = \sigma_1 + \cdots + \sigma_n > 0$ (car $A \neq 0$).

  \medskip
  \textit{Un chemin de rotations.}\enspace%
  On regroupe les coordonnées par paires consécutives $(1,2),(3,4),\ldots$ et,
  sur chaque paire, on place le bloc de rotation
  $\rho(\theta) = \left(\begin{smallmatrix}\cos\theta & -\sin\theta\\ \sin\theta & \cos\theta\end{smallmatrix}\right)$.
  \begin{itemize}
    \item si $n$ est \textbf{pair}, les $n/2$ paires couvrent tout : $R(\theta)$
          est diagonale par blocs $\rho(\theta)$, et $R(\theta)_{ii} = \cos\theta$
          pour tout $i$ ;
    \item si $n$ est \textbf{impair}, on laisse fixe la dernière coordonnée
          (celle du plus petit $\sigma_n$) et l'on fait tourner les $n-1$ autres.
  \end{itemize}
  Dans les deux cas $R(\theta) \in \mathrm{SO}_n(\mathbb{R}) \subset \mathrm{O}_n(\mathbb{R})$
  et $\theta \mapsto R(\theta)$ est continue. On calcule
  \[
    g(\theta) = \mathrm{tr}\bigl(S\,R(\theta)\bigr) =
    \begin{cases}
      \tau\cos\theta & (n \text{ pair}),\\[4pt]
      (\tau - \sigma_n)\cos\theta + \sigma_n & (n \text{ impair}).
    \end{cases}
  \]

  \medskip
  \textit{Conclusion par le théorème des valeurs intermédiaires.}\enspace%
  \index{théorème des valeurs intermédiaires}
  On a $g(0) = \tau > 0$, tandis que
  \[
    g(\pi) =
    \begin{cases}
      -\tau < 0 & (n \text{ pair}),\\[4pt]
      2\sigma_n - \tau \leq 0 & (n \text{ impair}),
    \end{cases}
  \]
  la dernière inégalité car $\sigma_n = \min_i \sigma_i \leq \tau/n \leq \tau/2$.
  La fonction continue $g$ change de signe (au sens large) sur $[0,\pi]$ : il
  existe $\theta_0 \in (0,\pi]$ tel que $g(\theta_0) = 0$. La matrice
  $R = R(\theta_0)$ répond à la question. \hfill$\blacksquare$
\end{mdframed}

\rubriq{Variante : argument de moyennisation}
On peut éviter la SVD au profit d'un argument de symétrie. Sur le groupe
compact connexe $\mathrm{SO}_n(\mathbb{R})$ ($n\geq 2$), la mesure de Haar
$\mathrm{d}Q$ vérifie $\int_{\mathrm{SO}_n} Q\,\mathrm{d}Q = 0$ : en effet
$M := \int Q\,\mathrm{d}Q$ satisfait $RM = M$ pour tout $R\in\mathrm{SO}_n$, et
le seul vecteur fixé par toute rotation est nul ($n\geq 2$). Dès lors
\[
  \int_{\mathrm{SO}_n} \mathrm{tr}(AQ)\,\mathrm{d}Q = \mathrm{tr}\!\Bigl(A\!\int Q\,\mathrm{d}Q\Bigr) = 0.
\]
La fonction continue $Q\mapsto\mathrm{tr}(AQ)$ est de moyenne nulle sur le
connexe $\mathrm{SO}_n$ : si elle ne s'annulait pas, elle garderait un signe
constant et son intégrale serait non nulle. Elle s'annule donc en un point.

\decsep{}

%=============================================================
%  § — EXEMPLES, REMARQUES, EXERCICES
%=============================================================
\secnum{exgreen}{4}{Exemples, remarques et exercices}

\noindent{\bfseries\color{navyblue} Exemples.}

\begin{mdframed}[style=exbar]
  \textbf{1. L'hyperplan de trace nulle ($A = I_n$).}\enspace%
  $H = \{M : \mathrm{tr}(M) = 0\}$ est le célèbre hyperplan
  $\mathfrak{sl}_n(\mathbb{R})$. La matrice de permutation cyclique $C$, de
  trace nulle et inversible, y vit : $C \in H \cap \mathrm{GL}_n$. Pour $n=2$,
  $C = \left(\begin{smallmatrix}0&1\\1&0\end{smallmatrix}\right)$ est même orthogonale,
  donc $C \in H \cap \mathrm{O}_2$.

  \smallskip\noindent
  \textbf{2. Cas $n=2$, $A = I_2$, version orthogonale fine.}\enspace%
  La rotation $\rho(\theta)$ a pour trace $2\cos\theta$, nulle pour
  $\theta = \pi/2$ : $\rho(\pi/2) = \left(\begin{smallmatrix}0&-1\\1&0\end{smallmatrix}\right)
  \in \mathrm{SO}_2 \cap H$. On illustre directement la Partie II.

  \smallskip\noindent
  \textbf{3. Hyperplan « coin supérieur gauche » ($A = E_{11}$, $r=1$).}\enspace%
  $H = \{M : M_{11} = 0\}$. La matrice de permutation cyclique a son coefficient
  $(1,1)$ nul, donc $C \in H \cap \mathrm{O}_n$ : ici une seule construction
  règle les deux parties.
\end{mdframed}

\vspace{6pt}
\noindent{\bfseries\color{counterred} Remarques et pièges.}

\begin{mdframed}[style=cexbar]
  \textbf{$n=1$ est exclu.}\enspace%
  L'unique hyperplan de $\mathcal{M}_1(\mathbb{R}) = \mathbb{R}$ est $\{0\}$,
  qui ne contient aucun scalaire non nul : ni $\mathrm{GL}_1$ ni $\mathrm{O}_1$
  ne sont rencontrés. L'hypothèse $n\geq 2$ est essentielle (elle sert au moment
  où l'on affirme que la diagonale de $C$ est nulle, ou qu'aucun vecteur non nul
  n'est fixé par toute rotation).

  \smallskip\noindent
  \textbf{« Rencontrer » n'est pas « contenir ».}\enspace%
  Un hyperplan ne contient jamais $\mathrm{GL}_n$ tout entier (qui engendre
  $\mathcal{M}_n$) ; le théorème dit seulement qu'il en contient \emph{au moins}
  un élément. De même il ne contient pas tout $\mathrm{O}_n$.

  \smallskip\noindent
  \textbf{La SVD n'est pas un luxe.}\enspace%
  Sans réduction, $\mathrm{tr}(AR(\theta))$ n'est pas une simple fonction de
  $\cos\theta$ : c'est la diagonalisation orthogonale des « directions » de $A$
  qui rend le calcul de $g(\theta)$ explicite.
\end{mdframed}

\vspace{6pt}
\exolabel

\begin{mdframed}[style=exobar]
  \begin{enumerate}
    \item Démontrer en détail que $M\mapsto\mathrm{tr}(AM)$ décrit \emph{toutes}
          les formes linéaires de $\mathcal{M}_n(\mathbb{R})$, et que $A$ est unique.

    \item Calculer $\det C$ pour la matrice de permutation cyclique
          $\gamma=(1\,2\,\cdots\,n)$ et vérifier $\mathrm{tr}(C)=0$ pour $n\geq 2$.

    \item Pour $n=3$ et $A = \mathrm{diag}(1,1,1)$ (hyperplan de trace nulle),
          exhiber explicitement une matrice de $\mathrm{O}_3 \cap H$ à l'aide du
          chemin de rotations de la Partie II. (\textit{Indication :} faire
          tourner deux coordonnées et garder la troisième fixe ;
          $g(\theta) = 2\cos\theta + 1$ s'annule en $\theta = 2\pi/3$.)

    \item Montrer que $\mathrm{Vect}\,\mathrm{GL}_n(\mathbb{R}) = \mathcal{M}_n(\mathbb{R})$
          et en déduire qu'aucun hyperplan ne contient $\mathrm{GL}_n$ tout entier.
          Comparer avec le théorème de la Partie I.

    \item \textbf{(Renforcement)} Montrer que, pour $n\geq 2$, tout hyperplan de
          $\mathcal{M}_n(\mathbb{R})$ rencontre $\mathrm{SL}_n(\mathbb{R})$.
          (\textit{Indication :} reprendre $N$ inversible de la Partie I avec
          $\sum_{i\leq r} N_{ii}=0$, puis ajuster une colonne pour normaliser le
          déterminant sans toucher aux coefficients diagonaux utiles.)

    \item \textbf{(Moyennisation)} Justifier $\int_{\mathrm{O}_n} Q\,\mathrm{d}Q = 0$
          pour $n\geq 2$ et en redéduire la Partie II sans recourir à la SVD.

    \item Le résultat de la Partie II subsiste-t-il en remplaçant $\mathrm{O}_n(\mathbb{R})$
          par $\mathrm{SO}_n(\mathbb{R})$~? (\textit{Oui :} le chemin construit
          reste dans $\mathrm{SO}_n$.) Et par $\mathrm{O}_n \setminus \mathrm{SO}_n$~?
  \end{enumerate}
\end{mdframed}

\leconfooter{Forme linéaire $M\mapsto\mathrm{tr}(AM)$ — équivalence \& valeurs singulières — argument de connexité}
</content>
</invoke>
